\section{The website and the demo application}
\label{sec:web}

Two surfaces present the results and both read the same JSON bundle, so neither can disagree with
the other about a score. Neither recomputes anything and neither writes anything back.

Build the bundle first, then the site:

\begin{cmd}
python web\build_site_data.py
cd web\website
npm run build
npm run preview
\end{cmd}

\noindent The first command reads \texttt{game\_scores.parquet} and \texttt{config.yaml} and writes
one JSON file per game plus an index; re-run it after any re-run of stage 7 or the site will show
the previous numbers. The build writes 977 pages into
\texttt{web\textbackslash website\textbackslash dist}, and \texttt{npm run preview} serves them at
\url{http://localhost:4321}. The built site is plain HTML, so there is no server to cold-start and
it works offline.

\noindent The closing line of the build is the one to check: 977 pages, one per scored game plus
the fixed pages.

\shotpair{site_home}{The site as it opens. The corpus size and collection date are stated on the
page rather than left implicit.}{site_browse}{The browse view. Per-game pages reachable from here
show the two components, the attribution, and the limitations applying to that title.}

The Streamlit application presents the same bundle on one screen, which suits a walkthrough:

\begin{cmd}
streamlit run web\app\app.py
\end{cmd}

\shot{streamlit_run}{0.98}{The demo application, reading the same bundle as the static site.}

